\chapter{A first project}
\label{ch:first}

This chapter takes you from an empty project to a verified output tree. It uses
the example design that comes with the manual.

\section{Before you start}

You can run every command of this chapter yourself. Work in a copy of the example
folder, because one step overwrites a file of the design. The example folder is
\file{docs/manual/examples} in the repository. In a POSIX shell, copy it with
\cmd{cp -r docs/manual/examples /tmp/vet-examples}, then \cmd{cd /tmp/vet-examples}.
In PowerShell, use \texttt{Copy-Item -Recurse docs\textbackslash manual\textbackslash examples
\$env:TEMP\textbackslash vet-examples}, then \texttt{cd \$env:TEMP\textbackslash vet-examples}.
You need \cmd{vet} installed (chapter~\ref{ch:install}).

The transcripts below are one session, shown in parts. A test replays each part in a
fresh copy of the folder, so each part starts with the commands of the parts before it
that it needs (\cmd{vet init}, \cmd{cp}, \cmd{cd}). If you follow along in one copy,
skip the commands that you have already typed. Every part that enters the folder
\file{myproj} leaves it again with \cmd{cd ..}, so you are back in the example folder.

\section{The example design}

The design has four source files and a file list. The file list names the
sources, one per line.

\lstinputlisting[style=source]{examples/design/design.f}

The top module is \code{soc_top}. It has a reset synchroniser and a peripheral
block. The reset synchroniser is a hand-written module.

\lstinputlisting[style=source]{examples/design/rtl/rst_sync.sv}

The peripheral block is generated. It instantiates the counter \code{timer}
twice. \code{timer} is a hand-written module too.

\lstinputlisting[style=source]{examples/design/rtl/periph_wrap.sv}
\lstinputlisting[style=source]{examples/design/rtl/timer.sv}

The top module is generated. It ties the input \code{test_i} of the
synchroniser to the constant \code{1'b0}.

\lstinputlisting[style=source]{examples/design/rtl/soc_top.sv}

Suppose you want a new input port \code{test_en_i} on \code{soc_top}, and you
want it to drive \code{test_i} of the synchroniser instead of the constant. This
is the job of the first script.

\section{The file list and the input tree}

hdl-vet reads the design from a \emph{file list}, a text file of the kind that
simulators read. It understands these elements:

\begin{itemize}
\item a source file, one per line;
\item \code{+incdir+a+b} for include folders, and \code{+define+A=1+B} for defines;
\item \code{-f FILE}, a nested list whose file names start from the base of the top-level
  list, and \code{-F FILE}, a nested list whose file names start from its own folder;
\item comments, from \code{//} to the end of the line, and lines that start with \texttt{\#};
\item environment variables, as \texttt{\$VAR}, \texttt{\$\{VAR\}} or \texttt{\$(VAR)}.
  A variable that is not set is an error.
\end{itemize}

Any other option is an error, so that nothing is dropped without notice.

Relative file names in the file list start from the \term{base}. The base is the folder of
the top-level file list, unless you give \opt{--base}. Use \opt{--base} when your tool
flow writes file lists for the folder where the simulator runs.

The \term{input tree} is the deepest folder that contains every source file, every
include folder and the base. In the example it is the folder \file{design}. The output
tree has the same layout. If the file list lies outside the input tree, hdl-vet writes
it at the top of the output tree, and the tree is not made wider by it. A file list that
sits far from its sources can still make the tree wider than you expect, when the base
is far from the sources too. Then give \opt{--base} with the folder that the sources are
relative to.

\section{Make a project}

A \term{project} is a folder with a project file, \file{hdl-vet.toml}, and your
scripts. The project file names the design, the scripts and the output folder, so
that you can run hdl-vet with no arguments. Chapter~\ref{ch:project} describes it
in full. You make a project with \cmd{vet init}.

\cmd{vet init} needs the file list (\opt{--filelist}, short \opt{-f}) and, if the
design has more than one top-level module, the name of the top module
(\opt{--top}). The example design has one, but we name it anyway. Always name the top
module in your own work. Chapter~\ref{ch:signoff} shows why.

\lstinputlisting[style=console]{examples/console/init.txt}

Read the transcript from the top.

\begin{itemize}
\item \cmd{vet init} loads the design. If the design has errors, \cmd{vet init}
  refuses. Otherwise it writes five files into the new folder \file{myproj}, and
  prints one line for each:
  the project file, a starter script, a \file{README.md}, a \file{.gitignore} and
  a note in the folder \file{cells}, where you keep the Verilog files that your scripts
  add (chapter~\ref{ch:operations}).
\item The project file names the file list relative to itself:
  \code{../design/design.f}.
\item The starter script changes nothing. It names the top module, lists the
  instances of the top module in comments, and shows operations that you can
  copy. It is a valid script: if you run \cmd{vet apply --dry-run} right after
  \cmd{vet init}, the answer is \code{0 applied, 0 skipped}.
\end{itemize}

\section{Write the first script}

A script is a Python file with one function, \code{apply}. hdl-vet calls it with
one argument, the \code{vet} object, which you use to find modules and to call
operations (chapter~\ref{ch:scripts}). Here is the first script.

\lstinputlisting[style=script]{examples/scripts/first.py}

The script does three things.

\begin{enumerate}
\item \code{vet.module("soc_top")} finds the module \code{soc_top} and gives
  you a handle for it.
\item \code{top.add_port("test_en_i", "input")} adds an input port. The type
  defaults to \code{logic}.
\item \code{connect("test_i", "test_en_i", was="1'b0")} connects the port
  \code{test_i} of the instance \code{u_rst_sync} to the new signal. The
  argument \code{was} says what you expect to find there now. If the design shows
  something else, for example after the generator changed, the operation stops
  with an error and changes nothing. This protects you from editing a design
  that is not the one you wrote the script for.
\end{enumerate}

In the next transcripts, the command \cmd{cp} puts this script in place of the starter
script, \file{myproj/scripts/01_example.py}. In your own project you edit that file, or
you add scripts to the list \code{scripts} in the project file.

\section{Look before you write}

The command \cmd{vet apply --dry-run} runs the scripts and the sign-off, and
prints what would change, as a unified diff. It writes no output tree.

\lstinputlisting[style=console]{examples/console/dry_run.txt}

The first line of the output is \code{project: ...}. Whenever hdl-vet uses a project
file, it says so first, so that you always know which project runs. The diff follows.
The last line, \code{2 applied, 0 skipped}, says that two operations changed the text.

Look at the diff. There is one line that you did not ask for. It is the header
line, \code{// Modified by hdl-vet; scripts: ...}. hdl-vet puts it into every
file that it changes, directly after the comment block at the top of the file. It
names each script and the first 12 hexadecimal digits of its hash. Look
also at the new port. It is not aligned with the other ports. hdl-vet never
reformats existing text, and the text it inserts is plain.

\section{Write the output tree}

\cmd{vet apply} without \opt{--dry-run} writes the output tree to the folder that
the project file names (\file{out}).

\lstinputlisting[style=console]{examples/console/run.txt}

The output is the folder, the counts, and one line for each changed file. The folder is
shown with its full name, because it comes from the project file.

If the output folder exists and is not empty, hdl-vet refuses to replace it,
unless you give \opt{--force}. Chapter~\ref{ch:cli} has the details.

\section{Look at the output tree}

The output tree has the same layout as the input tree.

\lstinputlisting[style=console]{examples/console/tree.txt}

\begin{itemize}
\item Files that no operation changed are exact copies.
\item The file list is written again, in \file{design.f}. It has one source per
  line. Nested lists are flattened into it, and it also carries the \code{+incdir+}
  and \code{+define+} lines. Its file names are relative to itself.
\item \file{hdl-vet.json} records what happened: the version, the top module,
  the scripts with their hashes, and for every changed file the hash before and
  after. \code{files} lists every file of the tree except \file{hdl-vet.json}.
  \code{added} lists files that scripts brought in (chapter~\ref{ch:operations}).
  There is no time stamp and no file name of your machine in it.
\end{itemize}

\section{Run it again on a regenerated design}

Next week the generator runs again. It writes a new \file{soc_top.sv}. In the
example, the new file has one more output, \code{busy_o}. You run the same
script on the new files.

\lstinputlisting[style=console]{examples/console/regenerated.txt}

The second command in the transcript, \cmd{cp}, does the work of the generator: it copies
the file \file{regenerated/rtl/soc_top.sv} over the design. This is why you work in a
copy. The same script runs again and gives the same edits, now on the new file. The
output tree has the new port \code{busy_o} and your new port \code{test_en_i}.

Two more things in the transcript.

\begin{itemize}
\item The output tree is not a valid input. Its files carry the header line, and
  hdl-vet refuses to read a file that carries it: \emph{the input must be an
  unmodified tree}. Always start from the generator's files.
\item The last command repeats the third command. The output folder
  \file{out/second} exists, so hdl-vet stops. The message names the option
  \opt{--force}, which replaces the folder.
\end{itemize}

The transcript uses \opt{--no-project}, which tells hdl-vet to ignore a project
file. The example folder has none. It is better to be explicit when you write
commands that other people run.
